\begin{recipe}
[% 
    preparationtime = {\unit[35]{min}},
    bakingtime = {\unit[25--30]{min}},
    bakingtemperature = {\protect\bakingtemperature{
        topbottomheat=\unit[200]{\textcelcius},
        fanoven=\unit[180]{\textcelcius}
    }},
    portion = {\portion{2}},
    %source = {},
    calory = {940kcal | 31g Protein | 104g KH | 42g Fett},
    price = {3,35€},
    created = {11.06.2026},
    %rating = {1},
]
{Feta-Spitzpaprika-Pasta aus dem Ofen}

\graph{% pictures
    big=pic/hauptgericht/18_feta-spitzpaprika-pasta
}

\introduction{%
Ofenpasta mit cremigem Feta, süßer Spitzpaprika und Knoblauch. Sehr wenig aktives Kochen, aber am Ende trotzdem eine richtig cremige Sauce. Also quasi Baked-Feta-Pasta, nur mit mehr Paprika und weniger TikTok-Trauma.
}

\ingredients{
    250 g & Pasta\index{Getreide!Pasta}\\
    200 g & Feta\index{Milchprodukte \& Alternativen!Feta}\\
    4 & Spitzpaprika\index{Gemüse!Spitzpaprika}\\
    1 kleine & Zwiebel\\
    2 Zehen & Knoblauch\\
    2--3 EL & Olivenöl\\
    1 TL & Paprikapulver\\
    \nicefrac{1}{2} TL & Oregano oder Thymian\\
    & Chiliflocken nach Geschmack\\
    & Pfeffer\\
    Optional & etwas Zitronensaft\\
    & etwas Pastawasser
}

\preparation{%
    \step%
    Backofen auf 200 °C Ober-/Unterhitze oder 180 °C Umluft vorheizen.
    
    \step%
    Spitzpaprika in Ringe oder Streifen schneiden. Zwiebel grob schneiden und Knoblauch fein hacken.
    
    \step%
    Spitzpaprika, Zwiebel und Knoblauch in eine Auflaufform geben. Den Feta in die Mitte legen.
    
    \step%
    Olivenöl darübergeben und alles mit Paprikapulver, Oregano oder Thymian, Chili und Pfeffer würzen. Mit Salz vorsichtig sein, da Feta schon salzig ist.
    
    \step%
    Die Auflaufform 25--30 Minuten backen, bis die Spitzpaprika weich ist und der Feta cremig wird.
    
    \step%
    Währenddessen Pasta in Salzwasser kochen. Vor dem Abgießen etwas Pastawasser auffangen.
    
    \step%
    Die Auflaufform aus dem Ofen nehmen und Feta, Spitzpaprika, Zwiebel und Knoblauch mit einer Gabel zu einer cremigen Sauce zerdrücken.
    
    \step%
    Pasta direkt in die Form geben und alles vermengen. Nach Bedarf etwas Pastawasser zugeben, bis die Sauce schön cremig ist.
    
    \step%
    Optional mit etwas Zitronensaft abschmecken und direkt servieren.
}

\hint{
    Pastawasser nicht vergessen. Ohne das wird es eher Feta-Gemüse-Matsch, mit Pastawasser wird es Sauce.
}

\suggestion[Variante]{%
    Mit etwas frischer Petersilie oder Basilikum wird das Gericht am Ende noch frischer. Wer es schärfer mag, kann die Chiliflocken direkt mit in den Ofen geben.
}

\end{recipe}